\section{Conclusion}\label{sec:conclusion}

High-dimensional chance-constrained simulation optimization cannot be made
sample efficient by acquisition design alone when a tiny target design misses
the feasible region.  We addressed that bottleneck with a source-learned,
dimension-equivariant structural proposal atlas, a replaceable online backend,
and an independent terminal verifier.  The theoretical analysis first rules
out unconditional finite-atlas coverage and then gives an explicit geometric
condition under which coverage is independent of nominal policy dimension.

The experiments identify the atlas as the dominant source of feasible-basin
coverage.  At dimension 1,000 it succeeds across three backends and all 60
trials, while matched common-Sobol designs almost always fail.  Canonical
SAASBO improves objective quality after coverage, but is not essential to
feasibility and is not a methodological contribution.  The dimension-10,000
and external energy studies support transfer of low-frequency structure, with
the important qualification that source learning does not beat an obvious
structured energy control.  Cumulative heteroscedastic decomposition improves
variance calibration but not optimization in the registered causal test.

The broader lesson is that historical simulations should be evaluated as an
experimental-design resource: what policy geometry do they cover, what target
assumption makes that geometry transferable, and how much independent evidence
is required before deployment?  Separating those questions from the optimizer
produces stronger attribution, more honest cost accounting, and a method whose
limits are as explicit as its gains.
